% ============================================================
% Section 5: System Architecture
% ============================================================
\section{System Architecture}

\subsection{Layered Architectural Overview}
SafeHer adopts a modular, three-tier architecture designed for scalability, maintainability, and real-time responsiveness. The system is partitioned into three logical layers: \textbf{Presentation (Frontend)}, \textbf{Application \& Data (Backend)}, and \textbf{Intelligence (AI)}. This separation enables independent evolution of components while ensuring secure, low-latency communication.

\textit{Presentation Layer}: Flutter mobile application with sensor integration (GPS, IMU), local state management via Riverpod, and offline persistence using Hive/SQLite.

\textit{Application \& Data Layer}: Firebase Authentication, Cloud Firestore (NoSQL), Cloud Functions (serverless), Firebase Cloud Messaging (FCM), and Crashlytics for monitoring.

\textit{Intelligence Layer}: Google Gemini Pro API with structured prompt engineering, risk score normalization, and fallback rule engine.

Data flows from mobile sensors through Firebase backend services to Google Gemini for AI inference, with autonomous alerting routed back via FCM to trusted contacts.

\subsection{Component Interactions}

\subsubsection{Frontend--Backend Communication}
Authentication uses phone OTP via Firebase Auth with JWT tokens attached to all subsequent requests. Data synchronization employs bi-directional Firestore listeners for real-time location/alert updates with an offline-first design and local queue for pending writes. Cloud Functions are invoked via HTTPS calls for danger score computation and alert dispatch.

\subsubsection{Backend--AI Layer Communication}
Gemini API calls are routed through authenticated Cloud Functions using a \textit{Secure Proxy Pattern} to prevent client-side key exposure. Context vectors are batched where privacy-compliant to optimize API quota. AI outputs are parsed, validated against schema, and normalized before score integration.

\subsubsection{Alert Distribution Pipeline}
When the Danger Score exceeds the threshold: (1)~Cloud Function validates user preferences; (2)~alert payload is constructed with user ID, location, risk score, and timestamp; (3)~parallel dispatch via FCM and SMS fallback (Twilio); (4)~log to Firestore alerts collection; (5)~initiate live location stream with 15-minute TTL.

\subsection{Security \& Privacy Architecture}
\begin{itemize}[leftmargin=*]
  \item \textbf{Data Minimization}: Only essential derived features (not raw coordinates) are transmitted for AI processing.
  \item \textbf{Encryption}: Data in transit (TLS 1.3); data at rest (Firestore server-side encryption); optional end-to-end encryption for alert payloads.
  \item \textbf{Access Control}: Firestore Security Rules enforce user-level document ownership; Cloud Functions verify caller identity via Firebase Admin SDK.
  \item \textbf{Auditability}: All critical operations are logged with immutable timestamps for forensic review.
\end{itemize}
